\documentclass[12pt]{article}
\usepackage[utf8]{inputenc}
\usepackage[T1]{fontenc}
\usepackage[margin=1in]{geometry}
\usepackage{setspace}
\usepackage{xcolor}
\usepackage[colorlinks=true, urlcolor=blue!60!black, linkcolor=blue!60!black]{hyperref}
\setlength{\parindent}{0em}
\setlength{\parskip}{0.45em}
\renewcommand{\baselinestretch}{1.0}
\begin{document}
\begin{center}
{\Large \textit{Saenz v. Roe}, 526 U.S. 489 (1999)\par}
\vspace{0.3em}
{\normalsize Case Brief --- One-Page Edition\par}
\end{center}

\textbf{Memory Jogger.} The right to travel and the revival of the Privileges or Immunities Clause: no second-class welfare benefits for new residents (\href{https://www.law.cornell.edu/supremecourt/text/98-97}{526 U.S. 489}).

\textbf{Facts.} California limited families resident under a year to the welfare grant of their prior state, as federal law (PRWORA) expressly allowed. New residents sued for a class.

\textbf{Procedural History.} The district court enjoined the cap; the Ninth Circuit affirmed.

\textbf{Votes.} Affirmed, 7--2: Stevens for the Court; Rehnquist and Thomas dissented.

\textbf{Holding.} New permanent residents have a Fourteenth Amendment right to be treated like other citizens of their state; discrimination by length of residence gets review ``no less strict'' than strict scrutiny, and Congress cannot authorize it.

\textbf{Reasoning.} Even \textit{Slaughter-House} agreed a citizen may become a citizen of any state by residence. Fiscal savings cannot justify ``degrees of citizenship based on length of residence,'' and welfare, unlike tuition, is not portable. Rehnquist said a settled resident is ``no longer `traveling' in any sense of the word''; Thomas said the Clause protects fundamental rights, not welfare.

\textbf{Examples.} \emph{Same-Side:} capping newcomers' housing stipends. \emph{Opposite-Side:} higher elk-license fees for visitors. \emph{Fence-Sitter:} a two-year tuition wait for settled adults.

\textbf{Critique.} Originalists split over a revival that ignored history; Tribe warned a citizens-only clause may weaken protection for noncitizens.

\textbf{Key Quotation.} ``The `right to travel' discussed in our cases embraces at least three different components. It protects the right of a citizen of one State to enter and to leave another State, the right to be treated as a welcome visitor rather than an unfriendly alien when temporarily present in the second State, and, for those travelers who elect to become permanent residents, the right to be treated like other citizens of that State.'' --- Justice Stevens, Opinion of the Court, 526 U.S. at 500.

\end{document}
